\subsection*{SRV Trace 9: The Conjugate Projection --- De-Operatorization and the Chromatic Solid}
\label{subsec:appB_srv_trace9}

\begin{quote}
\emph{``The wheel is not an illustration of the equations.
It is the equations with their operators removed.''}
\end{quote}

\noindent
Trace 8 (\S\ref{subsec:appB_srv_trace8}) factored the core dynamics in one
direction: parameters removed, operators retained, yielding the
\hyperref[sec:bk1_operatio]{Operatio} as the pre-parametric skeleton. This
trace performs the conjugate factoring. Remove the \emph{operators} ---
every act, map, flow, and derivation --- and ask what residue remains.
The answer is already in the manuscript, and has been since Book~V drew it:
the full chromatic coordinate of the Event Horizon Wheel
(Def.~\ref{definition:bk4_event_horizon_wheel},
Def.~\ref{definition:bk5_symbolic_shade}), rendered as the disk of
Fig.~\ref{figure:bk5_phi_chromatic_synthesis}. In categorical terms: Trace 8
forgot the objects and kept the arrows; Trace 9 forgets the arrows and keeps
the objects.

\subsubsection*{9.1 De-Operatorization: Equations $\to$ the Wheel}

We take the same core constructions as Trace 8 (\S8.1) and remove, from each,
everything that \emph{acts} --- differential operators, transport maps, flows,
traces, norms-as-constructions --- retaining only the datum each act operates
upon:

\begin{center}
\renewcommand{\arraystretch}{1.3}
\begin{tabularx}{\textwidth}{@{}lYY@{}}
\toprule
\textbf{Source} & \textbf{Operators removed} & \textbf{Parametric residue} \\
\midrule
Fokker--Planck (Thm.~\ref{theorem:bk1_fundamental_relation_fokker_plank_equation})
& $\partial/\partial s$, $\nabla\!\cdot$, $\nabla^2$, $D$ as act
& $\rho$: where the mass sits \\
Hamiltonian (Def.~\ref{definition:bk2_symbolic_hamiltonian})
& $\|\cdot\|_g$, $\mathrm{tr}$, the construction
& $H(x)$: how high, here \\
Equilibrium (Thm.~\ref{theorem:bk2_equilibrium_distribution})
& the equilibration itself
& $\beta$: how mixed \\
Transported overlap (Def.~\ref{definition:bk4_imaginary_symbolic_distance})
& transport, connection, flow $\Phi_s$
& $\Omega_O^\gamma = r\,e^{i\vartheta}$: phase and magnitude \\
\bottomrule
\end{tabularx}
\end{center}

\noindent
The last row is the canonical one. The transported overlap is the fundamental
observer-relative datum of the framework, and its operator-free residue is
exactly the pair that Book~V names: $\vartheta$ becomes hue
(Cor.~\ref{corollary:bk4_chromatic_transference_of_wheel}) and $r$ becomes
shade (Def.~\ref{definition:bk5_symbolic_shade}). The remaining residues are
scalar fields displayable \emph{on} that solid. The disk of
Fig.~\ref{figure:bk5_phi_chromatic_synthesis} is therefore not decoration
appended to the formalism: it is the total space of the parameter projection,
as the Operatio is the total expression of the operator projection.

\subsubsection*{9.2 The Asymmetry: What Composes, and What Appears}

The two projections are not symmetric, and the asymmetry is structural.

\begin{quote}
\textbf{Observation (composition asymmetry).}
The operator projection retains composability: $\vdash$ admits chaining, so
de-parameterized content supports sequence, cycle, and recurrence. The
parameter projection retains none: a residue with its arrows removed admits
no composition law. Argument: composition presupposes a directed morphism;
de-operatorization removes every morphism; a category stripped of its arrows
is a bare collection of objects.
\end{quote}

\noindent
Two consequences, both visible in the manuscript's own architecture. First,
the operator projection of PS requires \emph{four} pieces arranged in a cycle
--- \hyperref[sec:bk1_operatio]{Operatio},
\hyperref[sec:bk5_integratio]{Integratio},
\hyperref[sec:bk6_temperatio]{Temperatio},
\hyperref[sec:bk1_executio]{Executio}: differentiate, integrate, sample, free
--- because arrows generate order, and order generates stages. The parameter
projection requires exactly \emph{one} figure, presented all at once, because
without arrows there is no before. Time lives entirely on the operator side
of the factorization; the wheel is timeless not by idealization but by
construction. Second, the registers are forced: the operator projection can
still be \emph{said} (assertion is itself an arrow), so it is poetry; the
parameter projection can only be \emph{shown}, so it is a painting.

\subsubsection*{9.3 Universality of the Carrier}

The wheel does not borrow its color from optics, nor its resonance from sound.
Hue is the phase of an observer-relative overlap; shade is the magnitude of
accumulated reciprocal memory (Def.~\ref{definition:bk5_symbolic_shade}). By
Prop.~\ref{proposition:bk5_shade_transfers}(i) and the Modal Transference
Theorem (Thm.~\ref{theorem:appC_modal_transference}), the full pair
$(\mathrm{hue},\sigma)$ transfers intact to \emph{any} carrier preserving the
metric and adjacency structure of the overlap. Light and sound are carriers of
this geometry: instances of the transference, not its ground. Any bounded relational
process that carries a complex overlap datum --- an oscillation, an exchange,
a sustained relation (Scholium~\ref{scholium:bk5_palette_of_a_relation}), an
inference returned to and deepened --- has a position on the wheel: an angle
naming which mode it occupies, a radius naming how much has been laid down in
it. The scope of this claim is exactly the scope of the transference
hypotheses: proven under the stated conditions on the transference map,
neither narrower (mere optics) nor broader (unconditioned universality).

\subsubsection*{9.4 Framing}

Traces 8 and 9 together exhibit a factorization. The axiomatic Books are the
full object; the Operatio quartet is its exact projection onto operators; the
chromatic solid is its exact projection onto parameters. Neither projection
is a summary, an ornament, or an aid to intuition: each is the framework with
one of its two constituents removed, and the framework is recoverable from
neither alone. This is the architecture the work has carried from the
beginning --- one structure, shown at every resolution the bounded reader can
occupy: equations for the reader who derives, poetry for the reader who
composes, colour for the reader who sees.

\begin{flushright}
\emph{What acts, composes.\\
What remains, appears.\\
The book is their product.}
\end{flushright}
